% Options for packages loaded elsewhere
\PassOptionsToPackage{unicode}{hyperref}
\PassOptionsToPackage{hyphens}{url}
\documentclass[
]{article}
\usepackage{xcolor}
\usepackage{amsmath,amssymb}
\setcounter{secnumdepth}{-\maxdimen} % remove section numbering
\usepackage{iftex}
\ifPDFTeX
  \usepackage[T1]{fontenc}
  \usepackage[utf8]{inputenc}
  \usepackage{textcomp} % provide euro and other symbols
\else % if luatex or xetex
  \usepackage{unicode-math} % this also loads fontspec
  \defaultfontfeatures{Scale=MatchLowercase}
  \defaultfontfeatures[\rmfamily]{Ligatures=TeX,Scale=1}
\fi
\usepackage{lmodern}
\ifPDFTeX\else
  % xetex/luatex font selection
\fi
% Use upquote if available, for straight quotes in verbatim environments
\IfFileExists{upquote.sty}{\usepackage{upquote}}{}
\IfFileExists{microtype.sty}{% use microtype if available
  \usepackage[]{microtype}
  \UseMicrotypeSet[protrusion]{basicmath} % disable protrusion for tt fonts
}{}
\makeatletter
\@ifundefined{KOMAClassName}{% if non-KOMA class
  \IfFileExists{parskip.sty}{%
    \usepackage{parskip}
  }{% else
    \setlength{\parindent}{0pt}
    \setlength{\parskip}{6pt plus 2pt minus 1pt}}
}{% if KOMA class
  \KOMAoptions{parskip=half}}
\makeatother
\usepackage{color}
\usepackage{fancyvrb}
\newcommand{\VerbBar}{|}
\newcommand{\VERB}{\Verb[commandchars=\\\{\}]}
\DefineVerbatimEnvironment{Highlighting}{Verbatim}{commandchars=\\\{\}}
% Add ',fontsize=\small' for more characters per line
\newenvironment{Shaded}{}{}
\newcommand{\AlertTok}[1]{\textcolor[rgb]{1.00,0.00,0.00}{\textbf{#1}}}
\newcommand{\AnnotationTok}[1]{\textcolor[rgb]{0.38,0.63,0.69}{\textbf{\textit{#1}}}}
\newcommand{\AttributeTok}[1]{\textcolor[rgb]{0.49,0.56,0.16}{#1}}
\newcommand{\BaseNTok}[1]{\textcolor[rgb]{0.25,0.63,0.44}{#1}}
\newcommand{\BuiltInTok}[1]{\textcolor[rgb]{0.00,0.50,0.00}{#1}}
\newcommand{\CharTok}[1]{\textcolor[rgb]{0.25,0.44,0.63}{#1}}
\newcommand{\CommentTok}[1]{\textcolor[rgb]{0.38,0.63,0.69}{\textit{#1}}}
\newcommand{\CommentVarTok}[1]{\textcolor[rgb]{0.38,0.63,0.69}{\textbf{\textit{#1}}}}
\newcommand{\ConstantTok}[1]{\textcolor[rgb]{0.53,0.00,0.00}{#1}}
\newcommand{\ControlFlowTok}[1]{\textcolor[rgb]{0.00,0.44,0.13}{\textbf{#1}}}
\newcommand{\DataTypeTok}[1]{\textcolor[rgb]{0.56,0.13,0.00}{#1}}
\newcommand{\DecValTok}[1]{\textcolor[rgb]{0.25,0.63,0.44}{#1}}
\newcommand{\DocumentationTok}[1]{\textcolor[rgb]{0.73,0.13,0.13}{\textit{#1}}}
\newcommand{\ErrorTok}[1]{\textcolor[rgb]{1.00,0.00,0.00}{\textbf{#1}}}
\newcommand{\ExtensionTok}[1]{#1}
\newcommand{\FloatTok}[1]{\textcolor[rgb]{0.25,0.63,0.44}{#1}}
\newcommand{\FunctionTok}[1]{\textcolor[rgb]{0.02,0.16,0.49}{#1}}
\newcommand{\ImportTok}[1]{\textcolor[rgb]{0.00,0.50,0.00}{\textbf{#1}}}
\newcommand{\InformationTok}[1]{\textcolor[rgb]{0.38,0.63,0.69}{\textbf{\textit{#1}}}}
\newcommand{\KeywordTok}[1]{\textcolor[rgb]{0.00,0.44,0.13}{\textbf{#1}}}
\newcommand{\NormalTok}[1]{#1}
\newcommand{\OperatorTok}[1]{\textcolor[rgb]{0.40,0.40,0.40}{#1}}
\newcommand{\OtherTok}[1]{\textcolor[rgb]{0.00,0.44,0.13}{#1}}
\newcommand{\PreprocessorTok}[1]{\textcolor[rgb]{0.74,0.48,0.00}{#1}}
\newcommand{\RegionMarkerTok}[1]{#1}
\newcommand{\SpecialCharTok}[1]{\textcolor[rgb]{0.25,0.44,0.63}{#1}}
\newcommand{\SpecialStringTok}[1]{\textcolor[rgb]{0.73,0.40,0.53}{#1}}
\newcommand{\StringTok}[1]{\textcolor[rgb]{0.25,0.44,0.63}{#1}}
\newcommand{\VariableTok}[1]{\textcolor[rgb]{0.10,0.09,0.49}{#1}}
\newcommand{\VerbatimStringTok}[1]{\textcolor[rgb]{0.25,0.44,0.63}{#1}}
\newcommand{\WarningTok}[1]{\textcolor[rgb]{0.38,0.63,0.69}{\textbf{\textit{#1}}}}
\ifLuaTeX
  \usepackage{luacolor}
  \usepackage[soul]{lua-ul}
\else
  \usepackage{soul}
\fi
\setlength{\emergencystretch}{3em} % prevent overfull lines
\providecommand{\tightlist}{%
  \setlength{\itemsep}{0pt}\setlength{\parskip}{0pt}}
\usepackage{bookmark}
\IfFileExists{xurl.sty}{\usepackage{xurl}}{} % add URL line breaks if available
\urlstyle{same}
\hypersetup{
  hidelinks,
  pdfcreator={LaTeX via pandoc}}

\author{}
\date{}

\begin{document}

\subsection{0. 前言}\label{0-ux524dux8a00}

\begin{quote}
OI-wiki
感觉说的并不通俗易懂啊\ldots\ldots 导致它式子我还得自己手推一遍才能理解它的原理。
\\
写一篇面向普及组选手的 Min\_25 筛吧！\\
（不懂或有误可以指出并私信，必关）
\end{quote}

在学之前，先膜拜 Min\_25 大佬！！！

\subsection{1. 问题}\label{1-ux95eeux9898}

首先抛出一个问题：给你一个\textbf{积性函数} \(f\)，让你求
\(\displaystyle\sum_{i=1}^nf(i)\) ，也就是 \(f(i)\) 的前缀和，怎么做？\\
杜教筛通过构造一个 \(g\) 函数然后用卷积式子解决。\\
复杂度是 \(O(n^{\frac{2}{3}})\) 。\\
而如果 \(f(p)\) （\(p\)
为质数）能快速计算并表达成若干个完全积性函数的和（通常为关于 \(p\)
的多项式），且 \(f(p^c)\) 也可以快速计算时，就可以使用 Min\_25 筛。\\
完全积性函数的定义是，

\[\forall a,b,f(a)f(b)=f(ab)\]

注意，\(a,b\) \textbf{不需要互质}！

\subsection{2. 辅助函数}\label{2-ux8f85ux52a9ux51fdux6570}

我们先定义，\(G(k,n)\) 为:\\
从 \(\color{red}2\) 到 \(n\) 遍历整数 \(i\)，对于所有满足以下条件的
\(i\), 对 \(f(i)\) 求和。\\
条件： \(i\) 的最小质因子大于等于第 \(k\) 个质数。\\
用数学语言表述，就是：

\[\sum_{i=2}^nf(i)[L(i)\ge P_k]\]

其中 \(L(i)\) 就是 \(i\) 的最小质因子，\(P_k\) 就是第 \(k\) 个质数。\\
用代码表述，就是：

\begin{Shaded}
\begin{Highlighting}[]
\DataTypeTok{int}\NormalTok{ ans}\OperatorTok{=}\DecValTok{0}\OperatorTok{;}
\ControlFlowTok{for}\OperatorTok{(}\DataTypeTok{int}\NormalTok{ i}\OperatorTok{=}\DecValTok{2}\OperatorTok{;}\NormalTok{i}\OperatorTok{\textless{}=}\NormalTok{n}\OperatorTok{;}\NormalTok{i}\OperatorTok{++)\{}
		\ControlFlowTok{if}\OperatorTok{(}\NormalTok{L}\OperatorTok{[}\NormalTok{i}\OperatorTok{]\textgreater{}=}\NormalTok{pri}\OperatorTok{[}\NormalTok{k}\OperatorTok{])}\NormalTok{ ans}\OperatorTok{+=}\NormalTok{f}\OperatorTok{(}\NormalTok{i}\OperatorTok{);}
\OperatorTok{\}}
\ControlFlowTok{return}\NormalTok{ ans}\OperatorTok{;}
\end{Highlighting}
\end{Shaded}

边界就是当 \(P_k>n\) 时 \(G(k,n)=0\)。

\subsection{3. 转化}\label{3-ux8f6cux5316}

求这个奇怪的东西有什么用？\\
先不管，看看怎么求。\\
首先，我们发现，只有 \(L(i)\ge P_k\) 时，\(f(i)\) 才对 \(G(k,n)\)
产生贡献，所以我们枚举大于等于 \(P_k\) 的所有质数作为 \(i\)
的最小质因子。\\
具体怎么枚举呢？\\
对于每一个 \(P_j\)，考虑 \(L(i)=P_j\)，那么就有两种情况：

\begin{enumerate}
\def\labelenumi{\arabic{enumi}.}
\item
  \(P_j\) 在 \(i\) 中的次数等于 \(1\)，也就是说，\(i\) 因数分解后为
  \(\displaystyle\prod_{i=1}^k p_i^{c_i}\)，则 \(p_1=P_j,c_1=1\)。
\item
  \(P_j\) 在 \(i\) 中的次数大于 \(1\)，同理，\(p_1=P_j,c_1>1\)。
\end{enumerate}

\subsubsection{第一种情况}\label{ux7b2cux4e00ux79cdux60c5ux51b5}

对于所有满足第一种情况的 \(i\)，它们的总贡献就是 \\
\(f(P_j)+\sum_{t=2}^{\lfloor\frac{n}{P_j}\rfloor}f(tP_j)[L(t)\ge P_{j+1}]\)
\\
为啥得到这个式子呢？\\
首先因为 \(i\) 是 \(j\)
的倍数，所以我们按照学莫反时推式子的原理，转而去枚举
\(t=\dfrac {i}{P_j}\)。\\
然后因为 \(j\) 是 \(i\) 的最小质因子，所以 \(t\) 的最小素因子一定大于
\(P_j\)，或 \(t=1\)。\\
所以这个式子的前半部分处理的是 \(t=1\) 的情况，就是 \(f(P_j)\)。（此时
\(i=P_j\)）\\
\(t\ne 1\) 时，前面说过 \(t\) 的最小素因子一定大于 \(P_j\)，即大于等于
\(P_{j+1}\)。\\
所以 \(L(t)\ge P_{j+1}\) 时，\(f(i)\)，也就是 \(f(tP_j)\)
，才对答案构成贡献。\\
又因为 \(P_j\) 在 \(i\) 中的次数等于 \(1\)，所以
\(\gcd(t,P_j)=1\)，又因为 \(f\) 积性，所以 \(f(tP_j)=f(t)f(P_j)\)。\\
所以原式变为：\\
\(f(P_j)+f(P_j)\sum_{t=2}^{\lfloor\frac{n}{P_j}\rfloor}f(t)[L(t)\ge P_{j+1}]\)
根据 \(G\) 的定义，就可以写成：

\[f(P_j)+f(P_j)G(j+1,\lfloor\frac{n}{P_j}\rfloor)\]

\textbf{注意！下面是关键点！}\\
对于前面的一项 \(f(P_j),P_j\le n\) 即可，但对于后面一项
\(f(P_j)G(j+1,\lfloor\frac{n}{P_j}\rfloor)\),，必须 \(P_j^2\le n\)，即
\(t\ge P_j\)，也就是说，当 \(P_j^2>n\) 时，后面那一项就是无用的。 \\
因为后面那一项我们要求 \(t\ne 1\)，又 \(j\) 为 \(i\)
的\textbf{最小}素因子，所以必有 \(t\ge  P_j\) 。

\subsubsection{\texorpdfstring{第二种情况
}{第二种情况 }}\label{ux7b2cux4e8cux79cdux60c5ux51b5}

对于所有满足第一种情况的 \(i\)，它们的总贡献就是 \\
\(\sum_{c\ge2,P_j^c\le n}(f(P_j^c)+\sum_{t=2}^{\lfloor\frac{n}{P_j^c}\rfloor}f(tP_j^c)[L(t)\ge P_{j+1}])\)
我们枚举 \(c\) 作为 \(P_j\) 在 \(i\) 中的出现次数。\\
其余和第一种情况类似，也是枚举 \(t=\dfrac{i}{P_j^c}\)。\\
然后因为 \(j\) 是 \(i\) 的最小质因子，所以 \(t\) 的最小素因子一定大于
\(P_j\)，或 \(t=1\)。\\
所以这个式子的前半部分处理的是 \(t=1\) 的情况，就是 \(f(P_j^c)\)。（此时
\(i=P_j^c\)）\\
\(t\ne 1\) 时，前面说过 \(t\) 的最小素因子一定大于 \(P_j\)，即大于等于
\(P_{j+1}\)。\\
所以 \(L(t)\ge P_{j+1}\) 时，\(f(i)\)，也就是 \(f(tP_j^c)\)
，才对答案构成贡献。

然后同理继续推式子： \\
\(\sum_{c\ge2,P_j^c\le n}\left(f(P_j^c)+f(P_j^c)\sum_{t=2}^{\lfloor\frac{n}{P_j^c}\rfloor}f(t)[L(t)\ge P_{j+1}])\right)\)\\
还是根据 \(G\) 的定义：\\
\(\sum_{c\ge2,P_j^c\le n}\left(f(P_j^c)+f(P_j^c)G(j+1,\lfloor\frac{n}{P_j^c}\rfloor)\right)\)

\subsubsection{贡献}\label{ux8d21ux732e}

那么，对于所有最小质因子为 \(P_j\) 的 \(i\)，\(f(i)\) 的贡献之和就是：

\[f(P_j)+f(P_j)G(j+1,\lfloor\frac{n}{P_j}\rfloor)+\sum_{c \ge2,P_j^c\le n}(f(P_j^c)+f(P_j^c)G(j+1,\lfloor\frac{n}{P_j^c}\rfloor)\]

此时，对于每一个大于等于 \(k\) 的 \(j\)，计入贡献，则：

$$
\begin{aligned}
G(k,n)=&\left(\sum_{j\ge k,P_j\le n}\left(f(P_j)+f(P_j)G(j+1,\lfloor\frac{n}{P_j}\rfloor)+\sum_{c \ge2,P_j^c\le n}(f(P_j^c)+f(P_j^c)G(j+1,\lfloor\frac{n}{P_j^c}\rfloor)\right)\right)\\
=&\left(\sum_{j\ge k,P_j^2\le n}\left(f(P_j)G(j+1,\lfloor\frac{n}{P_j}\rfloor)+\sum_{c \ge2,P_j^c\le n}(f(P_j^c)+f(P_j^c)G(j+1,\lfloor\frac{n}{P_j^c}\rfloor))\right)\right)+\sum_{j\ge k,P_j\le n}f(P_j)
\end{aligned}
$$

这里说一下这一步是怎么推的。\\
由定义，当 \(P_k>n\) 时，\(G(k,n)=0\)，所以对于
\(\lfloor\sqrt{n}\rfloor< P_j\le n\) 的 \(P_j\)，必有
\(G (j+1,\lfloor\frac{n}{P_j}\rfloor)=0\)，所以不计入答案，又
\(c\ge2\)，所以 \(\lfloor\sqrt{n}\rfloor< P_j\le n\) 时只有 \(f(P_j)\)
那一项会计入贡献，那干脆把它提出去单独处理好了。

开始化简式子，

\[=\left(\sum_{j\ge k,P_j^2\le n}\left(\sum_{c \ge1,P_j^c\le n}\left(f(P_j^c)+f(P_j^c)G(j+1,\lfloor\frac{n}{P_j^c}\rfloor)\right)\right)-f(P_j)\right)+\sum_{j\ge k,P_j\le n}f(P_j)\]

这一步就很初等，没什么好说的。\\
下一步还是很初等。

\[G(k,n)=\left(\sum_{j\ge k,P_j^2\le n}\sum_{c \ge1,P_j^c\le n}f(P_j^c)([c\ne 1]+G(j+1,\lfloor\frac{n}{P_j^c}\rfloor))\right)+\sum_{j\ge k,P_j\le n}f(P_j)\]

设 \(S(n)=\displaystyle\sum_{p\le n,p\in\mathbb{P}}f(p)\)。\\
则可以表示成：

\[G(k,n)=\left(\sum_{j\ge k,P_j^2\le n}\sum_{c \ge1,P_j^c\le n}f(P_j^c)([c\ne 1]+G(j+1,\lfloor\frac{n}{P_j^c}\rfloor))\right)+S(n)-S(P_{k-1})\]

终于我们推了这么久，得到了这个递推式子。\\
终止条件就是当 \(P_k>n\) 是，\(G(k,n)=0\)。\\
而刚才那么多的分析，在 OI-wiki 中就是一个直接的等式，一笔带过。\\
\st{所以我说 OI-wiki 不好懂。} \\
这里有一个艾佛森括号，很丑。\\
怎么把它干掉？\\
简单！\\
我们知道，对于枚举到的最后一个 \(c\)，此时
\(P_j^c\le n\)，\(P_j^{c+1}>n\)，则有
\(\lfloor\frac{n}{P_j^c}\rfloor<P_j<P_{j+1}\)。\\
所以根据定义，此时的 \(G(j+1,\lfloor\frac{n}{P_j^c}\rfloor)=0\)。\\
所以我们就知道 \(c\) 枚举到最后一个的时候贡献只有 \(f(P_j^c)\)。\\
于是我们可以发现，根据乘法分配律，\(f (P_j^c)[c\ne 1]\) 这一项在 \(c=1\)
时没有贡献。\\
而 \(f(P_j^c)G(j+1,\lfloor\frac{n}{P_j^c}\rfloor)\) 这一项在最后一个
\(c\) 时没有贡献。 \\
所以我们可以枚举时将
\(\displaystyle\sum_{c \ge1,P_j^c\le n}f(P_j^c)([c\ne 1]+G(j+1,\lfloor\frac{n}{P_j^c}\rfloor))\)
，\\
改为
\(\displaystyle\sum_{c \ge1,P_{j}^{c+1}\le n}\left(f(P_j^c)G(j+1,\lfloor\frac{n}{P_j^c}\rfloor)+f(P_{j+1}^{c})\right)\)。\\
这一步很重要，务必理解透彻！

如果您还不理解，举个例子：\\
\(\sum_{i=1}^n(a_i+b_i)-a_1-b_n=\sum_{i=1}^{n-1}(a_{i+1}+b_i)\)\\
如果您理解了，继续推：\\
于是有

\[G(k,n)=\left(\sum_{j\ge k,P_j^2\le n}\sum_{c \ge1,P_{j}^{c+1}\le n}\left(f(P_j^c)G(j+1,\lfloor\frac{n}{P_j^c}\rfloor)+f(P_{j}^{c+1})\right)\right)+S(n)-S(P_{k-1})\]

假设 \(S\) 我们已经会计算了，那么这不就是一个 DP 式子吗？\\
直接模拟原式即可。

\subsection{\texorpdfstring{4. 求出
\(S\)}{4. 求出 S}}\label{4-ux6c42ux51fa--ux1d446}

现在考虑怎么求 \(S\) 。\\
我们知道 Min\_25 筛处理的 \(f\) 一般满足这样的性质：\\
对于质数 \(p\)，有 \(f(p)\) 是 \(p\) 的一个低次多项式。\\
设该多项式为 ： \\
\(\sum_{i=1}^r a_ip^{c_i}\)\\
\(r\) 就是项数，\(c_i\) 严格递增。（废话）\\
由 \(S(n)=\displaystyle\sum_{p\le n,p\in\mathbb{P}}f(p)\)，\\
得到：\\
\(S(n)=\displaystyle\sum_{p\le n,p\in\mathbb{P}}\displaystyle\sum_{i=1}^r a_ip^{c_i}=\displaystyle\sum_{i=1}^r a_i\displaystyle\sum_{p\le n,p\in\mathbb{P}}p^{c_i}\)

现在考虑对于每个 \(p^{c_i}\) 求贡献。\\
设 \(E(k,n)\)
为：\(\displaystyle\sum_{i=2}^n [P_k<L(i) \,\,\text{or}\,\, i\in\mathbb{P}]i^s\)，也就是埃氏筛在筛掉
\(k\) 轮之后剩下的所有 \(i\) 的 \(s\) 次方的和。\\
你不会埃氏筛？\\
好吧，\textbf{既然是一篇面向普及组选手的博客}，我还是说一下吧。\\
埃氏筛就是说，先筛掉 \(2\) 的倍数，再筛掉 \(3\) 的倍数，再筛掉 \(5\)
的倍数，以此类推，直到筛完小于等于 \(\lfloor\sqrt{n}\rfloor\)
的最大质数的倍数时，\textbf{\(\textbf{n}\)
及以下}没被筛掉的数就都是质数。\\
注意这里的``倍数''都是大于 \(2\) 的倍数。\\
为什么是 \(\lfloor\sqrt{n}\rfloor\) 呢？\\
考虑如果有一个 \(n\) 以内有一个合数 \(A\) 大于
\(\lfloor\sqrt{n}\rfloor\)，这就代表着它必有一个因数小于
\(\lfloor\sqrt{n}\rfloor\)，所以它一定已经被筛掉了。

你忘了什么是 \(P_k,L(i)\)? \st{（你是真行）}\\
再提一遍， \(L(i)\) 就是 \(i\) 的最小质因子，\(P_k\) 就是第 \(k\)
个质数。

既然这样，我们就可以得出，对于每个 \(s=c_i\)，它对 \(S(n)\) 的贡献就是
\(E(l,n)\)。\\
其中 \(P_l\) 为小于等于 \(\lfloor\sqrt{n}\rfloor\) 的最大质数。\\
现在就需要求出 \(E(k,n)\)。\\
首先，根据定义，\(E(k-1,n)\) 就是埃氏筛在筛掉 \(k-1\) 轮之后剩下的所有
\(i\) 的 \(s\) 次方的和。\\
我们就需要考虑第 \(k\) 轮筛完之后让 \(E(k-1,n)\) 的值减去了多少。\\
我们知道第 \(k\) 轮一定筛去了 \(P_k\) 的倍数, 而且。\\
这个减去的量就是：（开始推式子）


$$
\begin{aligned}
=&\sum_{i=2}^{\lfloor\frac{n}{P_k}\rfloor} [P_{k-1}<L(i) \,\,\text{or}\,\, i\in\mathbb{P}](iP_k)^s\\
=&P_k^s\sum_{i=2}^{\lfloor\frac{n}{P_k}\rfloor} [P_{k-1}<L(i) \,\,\text{or}\,\, i\in\mathbb{P}]i^s\\
=&P_k^sE(k-1,\lfloor\frac{n}{P_k}\rfloor)\\
\end{aligned}
$$

你一定会觉得这个式子有问题，或是看不懂，对吧？\\
\textbf{这个式子真对吗？}\\
首先，这个式子的意思是：对于 \(n\) 以内所有的 \(P_k\) 的倍数 \(m\)，若
\(L(\dfrac{m}{P_k})\) \textbf{大于等于} \(P_k\) （这就代表 \(L(m)\)
也大于等于 \(P_k\)）或者是 \(m\) 是 \(P_k\) 的质数倍，那么答案加上
\((mP_k)^s\)，这就是 \(E(k,n)\) 相对于 \(E(k-1,n)\) 所减去的量。

乍一看，好像没问题啊，就是埃氏筛第 \(k\) 轮所筛出的数的 \(s\)
次方的和啊，仔细一想，不对劲，如果 \(m\) 是 \(P_k\)
的质数倍，且这个质数小于 \(P_k\) ，那么 \(m\)
应该在之前就被筛掉了，怎么会留到现在呢？

所以说这个式子减多了。\\
怎么把这个式子加回来呢？\\
我们发现，减多了的数都可以被表示成如下式子：\(P_vP_k\)，其中 \(v<k\)。\\
那我们就枚举 \(P_v\) 及以内的数：

\[\sum_{i=2}^{P_{k-1}}[P_{k-1}<L(i)\,\,\text{or}\,\,i\in\mathbb{P}](iP_k)^s\]

我们发现 \(P_{k-1}\) 不可能小于 \(L(i)\)，所以原式就是：

\[\sum_{i=2}^{P_{k-1}}[i\in\mathbb{P}](iP_k)^s\]

这就是被减多了的式子。\\
你发现它就是：\(P_k^sE(k-1,P_{k-1})\)。\\
所以就可以得到最终的递推式：

\[E(k,n)=E(k-1,n)-P_k^s(E(k-1,\lfloor\frac{n}{P_k}\rfloor)-E(k-1,P_{k-1}))\]

\textbf{但是这个式子还是不一定对。}\\
当 \(P_k>\sqrt{n}\) 时，第 \(k\) 遍埃氏筛是不会筛掉任何数的。\\
所以真正的式子是：

\[E(k,n)=E(k-1,n)-P_k^s[P_k^2\le n](E(k-1,\lfloor\frac{n}{P_k}\rfloor)-E(k-1,P_{k-1}))\]

\textbf{这个式子就是对的了。}

最后递推的终止条件是
\(E(0,x)=\displaystyle\sum_{i=2}^xi^s\)。（我们假定第 \(0\) 个质数是
\(1\)）\\
那我们可以求出 \(E\) 了，也就可以求出 \(S\) 了。\\
也就可以求出 \(G\) 了。\\
我们发现一个性质：\(\displaystyle\sum_{i=1}^nf(i)=G(1,n)+f(1)=G(1,n)+1\)。\\
于是我们就基本学会了 Min\_25 筛的大致步骤。

\subsection{\texorpdfstring{5. 复杂度分析
}{5. 复杂度分析 }}\label{5-ux590dux6742ux5ea6ux5206ux6790}

在这里我不说具体复杂度是如何计算的，我在此想补充一些具体计算时的性质。

\subsubsection{优化}\label{ux4f18ux5316}

重新回顾前面 \(G(k,n)\) 的递推式：

\[G(k,n)=\left(\sum_{j\ge k,P_j^2\le n}\sum_{c \ge1,P_{j}^{c+1}\le n}\left(f(P_j^c)G(j+1,\lfloor\frac{n}{P_j^c}\rfloor)+f(P_{j}^{c+1})\right)\right)+S(n)-S(P_{k-1})\]

因为只有 \(P_j^2\le n\) 时两个 \(\Sigma\) 才计入贡献，才会计算
\(G(j+1,\lfloor\dfrac{n}{P_j^c}\rfloor)\)，因此可以得出在所有\textbf{需要算出的}
\(G(A,B)\) 中，\(P_A\) 必然小于等于
\(\lfloor\sqrt{n}\rfloor\)，\(P_{A-1}\) 则必然小于
\(\lfloor\sqrt{n}\rfloor\)。\\
然后我们有一个引理：

\[\lfloor\frac{\lfloor\frac{n}{x}\rfloor}{y}\rfloor=\lfloor\frac{n}{xy}\rfloor\]

所以 \(B\) 的所有取值都是某个 \(\lfloor\dfrac{n}{\lambda}\rfloor\) 。\\
我们又有一个引理：\\
\(\forall\lambda\le n,\lfloor\dfrac{n}{\lambda}\rfloor\) 最多只有
\(2\sqrt{n}\)
种取值：\(1,2,\cdots,\lfloor\sqrt{n}\rfloor,\lfloor\dfrac{n}{\sqrt{n}-1}\rfloor,\cdots,\lfloor\dfrac{n}{2}\rfloor,n\)。\\
又因为每次计算递推式时 \(S\) 函数只需要在 \(P_{A-1},B\)
点的取值就够了，所以在整个过程中，\(S\) 所需的取值点就只有
\(1,2,\cdots,\lfloor\sqrt{n}\rfloor,\lfloor\dfrac{n}{\sqrt{n}-1}\rfloor,\cdots,\lfloor\dfrac{n}{2}\rfloor,n\)
这 \(O(\sqrt{n})\) 个点。\\
但是 \(n\) 太大，你可以选择
\texttt{unordered\_map}，但是常数更小的方法是离散化然后存。

\subsubsection{\texorpdfstring{关于时间复杂度中求和与积分的转换
}{关于时间复杂度中求和与积分的转换 }}\label{ux5173ux4e8eux65f6ux95f4ux590dux6742ux5ea6ux4e2dux6c42ux548cux4e0eux79efux5206ux7684ux8f6cux6362}

\textbf{如果您不会积分可以跳过这一步，对算法学习无影响。}\\
我们发现，在许多算法的复杂度推导中，会发现这么一步：\\
（注意假定 \(f\) 是连续，恒为正，且可积的）

\[O(\sum_{i=1}^nf(i))=O(\int_0^nf(x)\mathrm{d}x)\]

为什么这么做是正确的？\\
有的人会说，因为求和就是离散的积分啊。\\
确实这么说也是对的，但是证一下更严谨。\\
即证：

\[\Theta(\sum_{i=1}^nf(i))=\int_0^nf(x)\mathrm{d}x\]

设
\(\sigma(n)=\displaystyle\sum_{i=1}^nf(i),\rho(n)=\displaystyle\int_0^nf(x)\mathrm{d}x\)。\\
即证：

\[\exists N,c_1,c_2>0,\forall n\ge N,n\in\mathbb{N}^+,c_1\sigma(n)\le\rho(n)\le c_2\sigma(n)\]

不妨直接证明更强的性质 \(N=0\) 时即满足。 \\
发现：

$$
\begin{aligned}
&c_1\sigma(n)\le\rho(n)\le c_2\sigma(n)\\
\iff&c_1\sum_{i=1}^nf(i)\le\int_0^nf(x)\mathrm{d}x\le c_2\sum_{i=1}^nf(i)
\end{aligned}
$$

那么我们只需要证，对 \(\forall i\in\mathbb{N}^+,1\le i\le n\)，有
\(u_{i}f(i)\le\displaystyle\int_{i-1}^if(x)\mathrm{d}x\le v_if(i)\)。\\
然后 \(c_1\) 取 \(\displaystyle\min_{i=1}^n u_i\)，\(c_2\) 取
\(\displaystyle\max_{i=1}^n v_i\) 即可。\\
由积分中值定理以及 \(f\) 连续可知，存在 \(\xi\in[i-1,i]\)，使得
\(\displaystyle\int_{i-1}^if(x)\mathrm{d}x=f(\xi)\),\\
取
\(u_i=\dfrac{f(\xi)}{f(i)}-\varepsilon,v_i=\dfrac{f(\xi)}{f(i)}+\varepsilon\)
即可。\\
得证。

由此我们有 min\_25 筛的复杂度是 \(O(\dfrac{n^{\frac{3}{4}}}{\log n})\)
或 \(O(n^{1-\varepsilon})\)，想看具体推导就见 OI-wiki
或者大佬们的博客吧。

\subsection{6. 具体实现}\label{6-ux5177ux4f53ux5b9eux73b0}

这里就以模板为例。

\begin{Shaded}
\begin{Highlighting}[]
\PreprocessorTok{\#include }\ImportTok{\textless{}bits/stdc++.h\textgreater{}}
\PreprocessorTok{\#define sor}\OperatorTok{(}\NormalTok{i}\OperatorTok{,}\PreprocessorTok{ }\NormalTok{l}\OperatorTok{,}\PreprocessorTok{ }\NormalTok{r}\OperatorTok{)}\PreprocessorTok{ }\ControlFlowTok{for}\PreprocessorTok{ }\OperatorTok{(}\DataTypeTok{int}\PreprocessorTok{ }\NormalTok{i}\PreprocessorTok{ }\OperatorTok{=}\PreprocessorTok{ }\NormalTok{l}\OperatorTok{;}\PreprocessorTok{ }\NormalTok{i}\PreprocessorTok{ }\OperatorTok{\textless{}=}\PreprocessorTok{ }\NormalTok{r}\OperatorTok{;}\PreprocessorTok{ }\NormalTok{i}\OperatorTok{++)}
\PreprocessorTok{\#define int }\DataTypeTok{\_\_int128}
\KeywordTok{using} \KeywordTok{namespace}\NormalTok{ std}\OperatorTok{;}
\CommentTok{/*  快读部分，暂时不予展示  */}
\KeywordTok{namespace}\NormalTok{ Revitalize}
\OperatorTok{\{}
    \AttributeTok{const} \DataTypeTok{int}\NormalTok{ N }\OperatorTok{=} \FloatTok{2e6}\OperatorTok{+}\DecValTok{7}\OperatorTok{,}\NormalTok{ P }\OperatorTok{=} \FloatTok{1e9} \OperatorTok{+} \DecValTok{7}\OperatorTok{;}
    \DataTypeTok{int}\NormalTok{ n}\OperatorTok{,}\NormalTok{m}\OperatorTok{;}
    \DataTypeTok{int}\NormalTok{ Ex01}\OperatorTok{[}\NormalTok{N}\OperatorTok{],}\NormalTok{ Ex02}\OperatorTok{[}\NormalTok{N}\OperatorTok{];}
    \DataTypeTok{int}\NormalTok{ ids}\OperatorTok{[}\NormalTok{N}\OperatorTok{],}\NormalTok{ sqrs}\OperatorTok{[}\NormalTok{N}\OperatorTok{];}
 
    \DataTypeTok{bool}\NormalTok{ vis}\OperatorTok{[}\NormalTok{N}\OperatorTok{],}\NormalTok{ cuur}\OperatorTok{;}
    \DataTypeTok{int}\NormalTok{ pri}\OperatorTok{[}\NormalTok{N}\OperatorTok{],}\NormalTok{Cl}\OperatorTok{[}\NormalTok{N}\OperatorTok{],}\NormalTok{I}\OperatorTok{[}\NormalTok{N}\OperatorTok{],}\NormalTok{II}\OperatorTok{[}\NormalTok{N}\OperatorTok{];}
    \KeywordTok{inline} \DataTypeTok{void}\NormalTok{ LinearSieve}\OperatorTok{(}\DataTypeTok{int}\NormalTok{ lim}\OperatorTok{)}
    \OperatorTok{\{}
\NormalTok{        ids}\OperatorTok{[}\DecValTok{1}\OperatorTok{]} \OperatorTok{=}\NormalTok{ sqrs}\OperatorTok{[}\DecValTok{1}\OperatorTok{]} \OperatorTok{=} \DecValTok{1}\OperatorTok{;}
        \ControlFlowTok{for} \OperatorTok{(}\DataTypeTok{int}\NormalTok{ i }\OperatorTok{=} \DecValTok{2}\OperatorTok{;}\NormalTok{ i }\OperatorTok{\textless{}=}\NormalTok{ lim}\OperatorTok{;}\NormalTok{ i}\OperatorTok{++)}
        \OperatorTok{\{}
            \ControlFlowTok{if} \OperatorTok{(!}\NormalTok{vis}\OperatorTok{[}\NormalTok{i}\OperatorTok{])}
            \OperatorTok{\{}
\NormalTok{                pri}\OperatorTok{[++}\NormalTok{pri}\OperatorTok{[}\DecValTok{0}\OperatorTok{]]} \OperatorTok{=}\NormalTok{ i}\OperatorTok{;}
\NormalTok{                ids}\OperatorTok{[}\NormalTok{pri}\OperatorTok{[}\DecValTok{0}\OperatorTok{]]} \OperatorTok{=} \OperatorTok{(}\NormalTok{ids}\OperatorTok{[}\NormalTok{pri}\OperatorTok{[}\DecValTok{0}\OperatorTok{]} \OperatorTok{{-}} \DecValTok{1}\OperatorTok{]} \OperatorTok{+}\NormalTok{ i}\OperatorTok{)} \OperatorTok{\%}\NormalTok{ P}\OperatorTok{;}
\NormalTok{                sqrs}\OperatorTok{[}\NormalTok{pri}\OperatorTok{[}\DecValTok{0}\OperatorTok{]]} \OperatorTok{=} \OperatorTok{(}\NormalTok{sqrs}\OperatorTok{[}\NormalTok{pri}\OperatorTok{[}\DecValTok{0}\OperatorTok{]} \OperatorTok{{-}} \DecValTok{1}\OperatorTok{]} \OperatorTok{+}\NormalTok{ i }\OperatorTok{*}\NormalTok{ i}\OperatorTok{)} \OperatorTok{\%}\NormalTok{ P}\OperatorTok{;}
            \OperatorTok{\}}
            \ControlFlowTok{for} \OperatorTok{(}\DataTypeTok{int}\NormalTok{ j }\OperatorTok{=} \DecValTok{1}\OperatorTok{;}\NormalTok{ j }\OperatorTok{\textless{}=}\NormalTok{ pri}\OperatorTok{[}\DecValTok{0}\OperatorTok{]} \OperatorTok{\&\&}\NormalTok{ i }\OperatorTok{*}\NormalTok{ pri}\OperatorTok{[}\NormalTok{j}\OperatorTok{]} \OperatorTok{\textless{}=}\NormalTok{ lim}\OperatorTok{;}\NormalTok{ j}\OperatorTok{++)}
            \OperatorTok{\{}
\NormalTok{                vis}\OperatorTok{[}\NormalTok{i }\OperatorTok{*}\NormalTok{ pri}\OperatorTok{[}\NormalTok{j}\OperatorTok{]]} \OperatorTok{=} \DecValTok{1}\OperatorTok{;}
                \ControlFlowTok{if} \OperatorTok{(}\NormalTok{i }\OperatorTok{\%}\NormalTok{ pri}\OperatorTok{[}\NormalTok{j}\OperatorTok{]} \OperatorTok{==} \DecValTok{0}\OperatorTok{)}
                \OperatorTok{\{}
                    \ControlFlowTok{break}\OperatorTok{;}
                \OperatorTok{\}}
            \OperatorTok{\}}
        \OperatorTok{\}}
    \OperatorTok{\}}
    \KeywordTok{inline} \DataTypeTok{int}\NormalTok{ L}\OperatorTok{(}\DataTypeTok{int}\NormalTok{ x}\OperatorTok{)\{}
        \ControlFlowTok{return} \OperatorTok{(}\NormalTok{x }\OperatorTok{\textless{}=}\NormalTok{ sqrt}\OperatorTok{((}\DataTypeTok{double}\OperatorTok{)}\NormalTok{n}\OperatorTok{))} \OperatorTok{?} \OperatorTok{(}\NormalTok{I}\OperatorTok{[}\NormalTok{x}\OperatorTok{])} \OperatorTok{:} \OperatorTok{(}\NormalTok{II}\OperatorTok{[}\NormalTok{n}\OperatorTok{/}\NormalTok{x}\OperatorTok{]);}
    \OperatorTok{\}}
    \KeywordTok{inline} \DataTypeTok{void}\NormalTok{ Ex}\OperatorTok{(}\DataTypeTok{int}\NormalTok{ n}\OperatorTok{)}
    \OperatorTok{\{}
        \ControlFlowTok{for} \OperatorTok{(}\DataTypeTok{int}\NormalTok{ k }\OperatorTok{=} \DecValTok{1}\OperatorTok{;}\NormalTok{ pri}\OperatorTok{[}\NormalTok{k}\OperatorTok{]} \OperatorTok{*}\NormalTok{ pri}\OperatorTok{[}\NormalTok{k}\OperatorTok{]} \OperatorTok{\textless{}=}\NormalTok{ n}\OperatorTok{;}\NormalTok{ k}\OperatorTok{++)}
        \OperatorTok{\{}
            \ControlFlowTok{for} \OperatorTok{(}\DataTypeTok{int}\NormalTok{ i }\OperatorTok{=} \DecValTok{1}\OperatorTok{;}\NormalTok{ i }\OperatorTok{\textless{}=}\NormalTok{ m }\OperatorTok{\&\&}\NormalTok{ pri}\OperatorTok{[}\NormalTok{k}\OperatorTok{]} \OperatorTok{*}\NormalTok{ pri}\OperatorTok{[}\NormalTok{k}\OperatorTok{]} \OperatorTok{\textless{}=}\NormalTok{ Cl}\OperatorTok{[}\NormalTok{i}\OperatorTok{];}\NormalTok{i}\OperatorTok{++)}
            \OperatorTok{\{}
\NormalTok{                Ex01}\OperatorTok{[}\NormalTok{i}\OperatorTok{]} \OperatorTok{=} \OperatorTok{(}\NormalTok{Ex01}\OperatorTok{[}\NormalTok{i}\OperatorTok{]} \OperatorTok{{-}}\NormalTok{ pri}\OperatorTok{[}\NormalTok{k}\OperatorTok{]} \OperatorTok{*} \OperatorTok{(}\NormalTok{Ex01}\OperatorTok{[}\NormalTok{L}\OperatorTok{(}\NormalTok{Cl}\OperatorTok{[}\NormalTok{i}\OperatorTok{]} \OperatorTok{/}\NormalTok{ pri}\OperatorTok{[}\NormalTok{k}\OperatorTok{])]} \OperatorTok{{-}}\NormalTok{ ids}\OperatorTok{[}\NormalTok{k }\OperatorTok{{-}} \DecValTok{1}\OperatorTok{])} \OperatorTok{\%}\NormalTok{ P }\OperatorTok{+}\NormalTok{ P }\OperatorTok{+}\NormalTok{ P}\OperatorTok{)} \OperatorTok{\%}\NormalTok{ P}\OperatorTok{;}
\NormalTok{                Ex02}\OperatorTok{[}\NormalTok{i}\OperatorTok{]} \OperatorTok{=} \OperatorTok{(}\NormalTok{Ex02}\OperatorTok{[}\NormalTok{i}\OperatorTok{]} \OperatorTok{{-}}\NormalTok{ pri}\OperatorTok{[}\NormalTok{k}\OperatorTok{]} \OperatorTok{*}\NormalTok{ pri}\OperatorTok{[}\NormalTok{k}\OperatorTok{]} \OperatorTok{\%}\NormalTok{ P }\OperatorTok{*} \OperatorTok{(}\NormalTok{Ex02}\OperatorTok{[}\NormalTok{L}\OperatorTok{(}\NormalTok{Cl}\OperatorTok{[}\NormalTok{i}\OperatorTok{]} \OperatorTok{/}\NormalTok{ pri}\OperatorTok{[}\NormalTok{k}\OperatorTok{])]} \OperatorTok{{-}}\NormalTok{ sqrs}\OperatorTok{[}\NormalTok{k }\OperatorTok{{-}} \DecValTok{1}\OperatorTok{])} \OperatorTok{\%}\NormalTok{ P }\OperatorTok{+}\NormalTok{ P }\OperatorTok{+}\NormalTok{ P}\OperatorTok{)} \OperatorTok{\%}\NormalTok{ P}\OperatorTok{;}
            \OperatorTok{\}}
        \OperatorTok{\}}
    \OperatorTok{\}}
    \DataTypeTok{int}\NormalTok{ G}\OperatorTok{(}\DataTypeTok{int}\NormalTok{ k}\OperatorTok{,} \DataTypeTok{int}\NormalTok{ n}\OperatorTok{)}
    \OperatorTok{\{}
        \ControlFlowTok{if} \OperatorTok{(}\NormalTok{pri}\OperatorTok{[}\NormalTok{k}\OperatorTok{]} \OperatorTok{\textgreater{}}\NormalTok{ n}\OperatorTok{)}
        \OperatorTok{\{}
            \ControlFlowTok{return} \DecValTok{0}\OperatorTok{;}
        \OperatorTok{\}}
        \DataTypeTok{int}\NormalTok{ ans }\OperatorTok{=} \DecValTok{0}\OperatorTok{;}
        \ControlFlowTok{for} \OperatorTok{(}\DataTypeTok{int}\NormalTok{ j }\OperatorTok{=}\NormalTok{ k}\OperatorTok{;}\NormalTok{ pri}\OperatorTok{[}\NormalTok{j}\OperatorTok{]} \OperatorTok{*}\NormalTok{ pri}\OperatorTok{[}\NormalTok{j}\OperatorTok{]} \OperatorTok{\textless{}=}\NormalTok{ n}\OperatorTok{;}\NormalTok{ j}\OperatorTok{++)}
        \OperatorTok{\{}
            \ControlFlowTok{for} \OperatorTok{(}\DataTypeTok{int}\NormalTok{ Pjc }\OperatorTok{=}\NormalTok{ pri}\OperatorTok{[}\NormalTok{j}\OperatorTok{];}\NormalTok{ Pjc }\OperatorTok{*}\NormalTok{ pri}\OperatorTok{[}\NormalTok{j}\OperatorTok{]} \OperatorTok{\textless{}=}\NormalTok{ n}\OperatorTok{;}\NormalTok{ Pjc }\OperatorTok{*=}\NormalTok{ pri}\OperatorTok{[}\NormalTok{j}\OperatorTok{])}
            \OperatorTok{\{}
\NormalTok{                ans }\OperatorTok{+=}\NormalTok{ Pjc }\OperatorTok{\%}\NormalTok{ P }\OperatorTok{*} \OperatorTok{(}\NormalTok{Pjc }\OperatorTok{{-}} \DecValTok{1} \OperatorTok{+}\NormalTok{ P}\OperatorTok{)} \OperatorTok{\%}\NormalTok{ P }\OperatorTok{*}\NormalTok{ G}\OperatorTok{(}\NormalTok{j }\OperatorTok{+} \DecValTok{1}\OperatorTok{,}\NormalTok{ n }\OperatorTok{/}\NormalTok{ Pjc}\OperatorTok{)} \OperatorTok{\%}\NormalTok{ P}\OperatorTok{;}
\NormalTok{                ans }\OperatorTok{\%=}\NormalTok{ P}\OperatorTok{;}
\NormalTok{                ans }\OperatorTok{+=} \OperatorTok{(}\NormalTok{Pjc }\OperatorTok{\%}\NormalTok{ P }\OperatorTok{*}\NormalTok{ pri}\OperatorTok{[}\NormalTok{j}\OperatorTok{]} \OperatorTok{\%}\NormalTok{ P}\OperatorTok{)} \OperatorTok{*} \OperatorTok{(}\NormalTok{Pjc }\OperatorTok{*}\NormalTok{ pri}\OperatorTok{[}\NormalTok{j}\OperatorTok{]} \OperatorTok{\%}\NormalTok{ P }\OperatorTok{{-}} \DecValTok{1} \OperatorTok{+}\NormalTok{ P}\OperatorTok{)} \OperatorTok{\%}\NormalTok{ P}\OperatorTok{;}
\NormalTok{                ans }\OperatorTok{\%=}\NormalTok{ P}\OperatorTok{;}
            \OperatorTok{\}}
        \OperatorTok{\}}
\NormalTok{        ans }\OperatorTok{+=} \OperatorTok{(}\NormalTok{Ex02}\OperatorTok{[}\NormalTok{L}\OperatorTok{(}\NormalTok{n}\OperatorTok{)]} \OperatorTok{{-}}\NormalTok{ Ex01}\OperatorTok{[}\NormalTok{L}\OperatorTok{(}\NormalTok{n}\OperatorTok{)]} \OperatorTok{+}\NormalTok{ P}\OperatorTok{)} \OperatorTok{\%}\NormalTok{ P}\OperatorTok{;}
\NormalTok{        ans }\OperatorTok{{-}=} \OperatorTok{(}\NormalTok{sqrs}\OperatorTok{[}\NormalTok{k }\OperatorTok{{-}} \DecValTok{1}\OperatorTok{]} \OperatorTok{{-}}\NormalTok{ ids}\OperatorTok{[}\NormalTok{k }\OperatorTok{{-}} \DecValTok{1}\OperatorTok{]} \OperatorTok{+}\NormalTok{ P}\OperatorTok{)} \OperatorTok{\%}\NormalTok{ P}\OperatorTok{;}
        \ControlFlowTok{return}\NormalTok{ ans}\OperatorTok{;}
    \OperatorTok{\}}
    \KeywordTok{inline} \DataTypeTok{void}\NormalTok{ Lx}\OperatorTok{(}\DataTypeTok{int}\NormalTok{ n}\OperatorTok{)}
    \OperatorTok{\{}
        \DataTypeTok{int}\NormalTok{ l }\OperatorTok{=} \DecValTok{1}\OperatorTok{,}\NormalTok{ r }\OperatorTok{=} \DecValTok{0}\OperatorTok{;}
        \ControlFlowTok{while} \OperatorTok{(}\NormalTok{l}\OperatorTok{\textless{}=}\NormalTok{n}\OperatorTok{)\{}
\NormalTok{            r }\OperatorTok{=}\NormalTok{ n }\OperatorTok{/} \OperatorTok{(}\NormalTok{n }\OperatorTok{/}\NormalTok{ l}\OperatorTok{);}
\NormalTok{            Cl}\OperatorTok{[++}\NormalTok{m}\OperatorTok{]} \OperatorTok{=}\NormalTok{ n }\OperatorTok{/}\NormalTok{ l}\OperatorTok{;}
\NormalTok{            Ex01}\OperatorTok{[}\NormalTok{m}\OperatorTok{]} \OperatorTok{=}\NormalTok{ Cl}\OperatorTok{[}\NormalTok{m}\OperatorTok{]} \OperatorTok{*} \OperatorTok{(}\NormalTok{Cl}\OperatorTok{[}\NormalTok{m}\OperatorTok{]} \OperatorTok{+} \DecValTok{1}\OperatorTok{)} \OperatorTok{/} \DecValTok{2} \OperatorTok{\%}\NormalTok{ P }\OperatorTok{{-}} \DecValTok{1}\OperatorTok{;}
\NormalTok{            Ex02}\OperatorTok{[}\NormalTok{m}\OperatorTok{]} \OperatorTok{=}\NormalTok{ Cl}\OperatorTok{[}\NormalTok{m}\OperatorTok{]} \OperatorTok{*} \OperatorTok{(}\NormalTok{Cl}\OperatorTok{[}\NormalTok{m}\OperatorTok{]} \OperatorTok{+} \DecValTok{1}\OperatorTok{)} \OperatorTok{\%}\NormalTok{ P }\OperatorTok{*} \OperatorTok{(}\NormalTok{Cl}\OperatorTok{[}\NormalTok{m}\OperatorTok{]} \OperatorTok{+}\NormalTok{ Cl}\OperatorTok{[}\NormalTok{m}\OperatorTok{]} \OperatorTok{+} \DecValTok{1}\OperatorTok{)} \OperatorTok{\%}\NormalTok{ P }\OperatorTok{*} \DecValTok{166666668} \OperatorTok{\%}\NormalTok{ P }\OperatorTok{{-}} \DecValTok{1}\OperatorTok{;}
            \ControlFlowTok{if}\OperatorTok{(}\NormalTok{Cl}\OperatorTok{[}\NormalTok{m}\OperatorTok{]\textless{}=}\NormalTok{sqrt}\OperatorTok{((}\DataTypeTok{double}\OperatorTok{)(}\NormalTok{n}\OperatorTok{)))}
\NormalTok{                I}\OperatorTok{[}\NormalTok{Cl}\OperatorTok{[}\NormalTok{m}\OperatorTok{]]} \OperatorTok{=}\NormalTok{ m}\OperatorTok{;}
            \ControlFlowTok{else}
\NormalTok{                II}\OperatorTok{[}\NormalTok{n }\OperatorTok{/}\NormalTok{ Cl}\OperatorTok{[}\NormalTok{m}\OperatorTok{]]} \OperatorTok{=}\NormalTok{ m}\OperatorTok{;}
\NormalTok{            l }\OperatorTok{=}\NormalTok{ r }\OperatorTok{+} \DecValTok{1}\OperatorTok{;}
        \OperatorTok{\}}
    \OperatorTok{\}}
    \KeywordTok{inline} \DataTypeTok{void}\NormalTok{ work}\OperatorTok{()}
    \OperatorTok{\{}
\NormalTok{        read}\OperatorTok{(}\NormalTok{n}\OperatorTok{);}
\NormalTok{        LinearSieve}\OperatorTok{((}\DataTypeTok{int}\OperatorTok{)(}\DecValTok{1000000}\OperatorTok{));}
\NormalTok{        Lx}\OperatorTok{(}\NormalTok{n}\OperatorTok{);}
\NormalTok{        Ex}\OperatorTok{(}\NormalTok{n}\OperatorTok{);}
        \DataTypeTok{int}\NormalTok{ ans }\OperatorTok{=}\NormalTok{ G}\OperatorTok{(}\DecValTok{1}\OperatorTok{,}\NormalTok{ n}\OperatorTok{)} \OperatorTok{\%}\NormalTok{ P }\OperatorTok{+} \DecValTok{1}\OperatorTok{;}
\NormalTok{        print}\OperatorTok{(}\NormalTok{ans}\OperatorTok{);}
    \OperatorTok{\}}
\OperatorTok{\}}
\DataTypeTok{signed}\NormalTok{ main}\OperatorTok{()}
\OperatorTok{\{}
\NormalTok{    ios}\OperatorTok{::}\NormalTok{sync\_with\_stdio}\OperatorTok{(}\KeywordTok{false}\OperatorTok{);}
\NormalTok{    cin}\OperatorTok{.}\NormalTok{tie}\OperatorTok{(}\DecValTok{0}\OperatorTok{),}\NormalTok{ cout}\OperatorTok{.}\NormalTok{tie}\OperatorTok{(}\DecValTok{0}\OperatorTok{);}
    \ControlFlowTok{return}\NormalTok{ Revitalize}\OperatorTok{::}\NormalTok{work}\OperatorTok{(),} \DecValTok{0}\OperatorTok{;}
\OperatorTok{\}}
\end{Highlighting}
\end{Shaded}

顺便警示一下后人，这题需要 \texttt{\_\_int128}，而且要随时随地取模。

\subsection{7. 优化}\label{7-ux4f18ux5316}

这里我知道有三种优化：

\begin{itemize}
\item
  \href{https://www.luogu.com.cn/article/gw5rqtsc}{Min\_26
  筛}，使用根号分治优化。
\item
  \href{https://zhidao.baidu.com/question/692919259662897292.html}{新版
  Min\_25 筛}，然而公式炸了，似乎是树状数组优化？
\item
  \href{https://www.cnblogs.com/zkyJuruo/p/17544928.html}{zhoukangyang
  的算法}，优化到了 \(O(\sqrt{n}\mathrm{polylog}(n))\)，太强了！\\
  \st{这像是我能会的吗？}
\end{itemize}

\subsection{Inf. 鸣谢}\label{inf-ux9e23ux8c22}

感谢 \href{https://luogu.com.cn/user/557800}{luuia}
提供的高质量模拟赛使得我学习这一算法。\\
感谢 \href{https://luogu.com.cn/user/716122}{Collapsarr}
提出在赛时当场学当场写的新思路。（尽管没有实现）\\
感谢神的光环照耀。\\
感谢 CWH 提供精神支持。

\end{document}
